% engines: lualatex
% expect-warnings: 0
% expect-errors: 0
\documentclass{article}
\usepackage[margin=1cm]{geometry}
\usepackage[kernel = lua]{../../latex/pgf-go}
\newcommand\case[1]{\par\medskip\noindent\textbf{#1}\par\nopagebreak}
\pagestyle{empty}

% Draws frames of \sgfanimation one by one (as in sgf-animation.tex).
\ExplSyntaxOn
\NewDocumentCommand \frames {mmm} {
	\group_begin:
		\pgfgo_sgfanimation_setup:nn {#1} {#2}
		\clist_map_inline:nn {#3} {\pgfgo_sgfanimation_frame:n {##1}\quad}
	\group_end:
}
\ExplSyntaxOff
\begin{document}

% partial = <corners> in \sgfdiagrams, \sgfproblem and \sgfanimation: the same part in every board,
% also when the file has a view (VW).

\case{sgf diagrams, partial = K10:T19: every diagram, also the variation}
\sgfdiagrams[partial = K10:T19, step = 2, variations, goban = {scale = 0.5}]{joseki.sgf}

\case{sgf diagrams, partial = \{T19, N12\}}
\sgfdiagrams[partial = {T19, N12}, goban = {scale = 0.5}]{joseki.sgf}

\case{sgf diagrams: partial = A1:I9 wins over the views of the file (the whole board in every diagram)}
\sgfdiagrams[partial = A1:I9, step = 2, goban = {scale = 0.4}]{views.sgf}

\case{sgf problem, partial = A1:T19: the whole board}
\sgfproblem[solutions = false, partial = A1:T19, goban = {scale = 0.3}]{problem.sgf}

\case{sgf animation, partial = M12:T19: frames 0, 2 and 5}
\frames{partial = M12:T19, goban = {scale = 0.4}}{joseki.sgf}{0, 2, 5}

\end{document}
